\section{Correspondence Map: This Paper \texorpdfstring{$\leftrightarrow$}{<->} Conceptual Topography}


For readers who will also read Conceptual Topography (CT) alongside this paper, the correspondence between the main concepts of both papers is shown here.

This paper and CT describe different projected windows of the same structure of differential development. The variables of the two papers stand in correspondence, but they are not simply identical across layers. Neither paper grounds the other. From different observation windows, the two papers describe the state that opens as consciousness and the terrain that stabilizes as concepts as different projections of the same generative grammar (VED×IFGT). The projected-uses stance fixed in §\ref{sec:projected-uses-stance} runs through the whole of this correspondence map.

\begin{table}[H]
\small
\begin{tabularx}{\textwidth}{>{\raggedright\arraybackslash}X >{\raggedright\arraybackslash}X >{\raggedright\arraybackslash}X}
\toprule
This paper & CT (Conceptual Topography) & Shared variables / correspondence \\
\midrule
Slow-layer sedimentation structure & Information potential $\Phi$ landscape & Both converge as $\Phi^*$, the long-time attractor \\
SSO (temporal non-closure) & Invisible GPU (landscape operating below consciousness) & $\Phi$ constrains thought without being directly referenced by consciousness \\
TCA (transformation of speed difference) & Token coupling $\lambda\sum_i \ell_i \cdot \nabla\Phi$ & Isomorphic structures transforming speed/spatial difference \\
Log-referencing direction by meaning gradient $\nabla\Phi$ & Concept definition: locally stable region $\nabla\Phi\approx 0$ & $\nabla\Phi$ simultaneously governs conscious bias and concept formation \\
Log referencing (state where consciousness opens) & Traversal cost reduction (state where understanding obtains) & Both described as referencing/traversal of $\nabla\Phi$ structure \\
Qualia (live foreground at sedimentation boundary) & Insight (rapid topological transition) & The same boundary phenomenon observed at different temporal granularities \\
$I(x,\tau) > 0$ (exclusion of complete closure) & $I(x,\tau) > 0$ (exclusion of complete conceptual closure) & Positivity marks the exclusion of ordinary complete closure in each window; it does not by itself identify consciousness with the conceptual landscape \\
\bottomrule
\end{tabularx}
\end{table}

\textbf{Reading suggestions:}

For readers who read this paper first: reading CT reveals that the "temporal condition for consciousness to open" in this paper and the "spatial condition for concepts to form" in CT come from the same structure. What is referenced in a conscious state is determined by the geometry of the conceptual landscape the person holds.

For readers who read CT first: reading this paper reveals that behind CT's description of "thought flowing across terrain," the speed difference between fast and slow layers is required. The terrain is not given; it becomes referenceable only when SSO is established.

Read in mutual reference, the two papers suggest the following. The state that opens as consciousness and the terrain that stabilizes as concepts are different projected stabilization regimes of the same VED×IFGT grammar, not identical entities or reducible objects. The same policy extends to the VED physical layer, the Bio-IFGT biological layer, and the intelligence layer treated in the intelligence paper \cite{Intelligence}; the four layers as a whole are read as different projected windows of the same differential-development structure. The consciousness window of this paper is one of those windows, and its correspondence with the other windows should be understood not as reduction but as structural isomorphism under the effective coordinates proper to each window.
